% ============================ CHAPITRE 2 ==============================
\chapter{Analyse et spécification des besoins}

\section*{Introduction}
Ce chapitre identifie les acteurs du système, spécifie les besoins fonctionnels
et non fonctionnels, puis les illustre au moyen du diagramme de cas
d'utilisation et de la description de quelques cas représentatifs.

\section{Identification des acteurs}
Trois profils interagissent avec l'application, conformément aux rôles gérés par
le module de sécurité (table \texttt{SECURITY\_USERS}). Le
tableau~\ref{tab:acteurs} en précise les responsabilités.

\begin{table}[htbp]
\centering
\caption{Identification des acteurs}
\label{tab:acteurs}
\renewcommand{\arraystretch}{1.25}
\begin{tabularx}{\textwidth}{>{\bfseries}p{4cm}X}
\toprule
\textmd{\textbf{Acteur}} & \textbf{Responsabilités}\\
\midrule
Administrateur (\texttt{ADMIN}) & Gère les agences, les employés et les clients ;
valide les demandes de création de clients et d'employés ; consulte le journal
d'audit.\\
Directeur commercial (\texttt{DIRECTEUR\_COMMERCIAL}) & Soumet les demandes de
création (clients, employés) ; valide les demandes de modification des données
des employés ; consulte les objectifs.\\
Utilisateur (\texttt{USER}) & Consulte les données autorisées et soumet des
demandes de modification de ses propres informations.\\
\bottomrule
\end{tabularx}
\end{table}

\section{Besoins fonctionnels}
Le système doit permettre :
\begin{itemize}[noitemsep]
  \item l'authentification sécurisée et la gestion des rôles et des droits d'accès ;
  \item la gestion (CRUD) des agences, des employés et des clients ;
  \item le rattachement des employés et des clients à leur agence et à leur
        gestionnaire ;
  \item la consultation des objectifs commerciaux et le suivi de la production ;
  \item le \textbf{pointage (présence) des employés} : saisie de l'heure d'arrivée
        et de départ et détermination du statut (présent / retard / absent),
        avec saisie automatique (bouton) ou manuelle (administrateur) ;
  \item la soumission et la validation des demandes (clients, employés,
        modifications) selon le cycle \texttt{EN\_ATTENTE} $\rightarrow$
        \texttt{EN\_COURS} $\rightarrow$ \texttt{EFFECTUE}/\texttt{REFUSE} ;
  \item la notification des demandes en attente selon le rôle ;
  \item la journalisation des actions (audit) et l'affichage de l'activité
        récente ;
  \item la restitution d'indicateurs via un tableau de bord et via Power BI.
\end{itemize}

\section{Besoins non fonctionnels}
Le tableau~\ref{tab:bnf} récapitule les besoins non fonctionnels.

\begin{table}[htbp]
\centering
\caption{Besoins non fonctionnels}
\label{tab:bnf}
\renewcommand{\arraystretch}{1.25}
\begin{tabularx}{\textwidth}{>{\bfseries}p{3.4cm}X}
\toprule
\textmd{\textbf{Besoin}} & \textbf{Description}\\
\midrule
Sécurité & Authentification, gestion des rôles et protection des opérations
sensibles par des workflows de validation.\\
Fiabilité & Intégrité des données garantie par le SGBD et les transactions JTA.\\
Performance & Temps de réponse acceptable sur les volumes manipulés.\\
Ergonomie & Interface claire et cohérente, navigation par menu latéral.\\
Maintenabilité & Architecture en couches, code modulaire et faiblement couplé.\\
Portabilité & Application Jakarta EE standard, déployable sur serveur compatible.\\
\bottomrule
\end{tabularx}
\end{table}

\section{Diagramme de cas d'utilisation global}
La figure~\ref{fig:uc_global} présente le \textbf{diagramme de cas d'utilisation
global}. Ce diagramme UML a pour objectif de représenter, d'un point de vue
fonctionnel, les interactions entre les acteurs et les services offerts par le
système.

\begin{figure}[htbp]
  \centering
  \includegraphics[width=0.95\textwidth]{uc_global.png}
  \caption{Diagramme de cas d'utilisation global}
  \label{fig:uc_global}
\end{figure}

On y distingue les trois acteurs identifiés et leurs cas d'utilisation
respectifs. Le cas \emph{S'authentifier} est commun à tous les acteurs ; il est
relié aux autres cas par une relation \texttt{\guillemotleft include\guillemotright}
traduisant le fait que l'accès aux fonctionnalités exige une authentification
préalable. L'\textbf{administrateur} concentre les cas de gestion (agences,
employés, clients), la validation des demandes et la consultation du journal ; le
\textbf{directeur commercial} soumet les demandes et valide les modifications ;
l'\textbf{utilisateur} se limite à la consultation et à la soumission de demandes
le concernant.

\section{Description de quelques cas d'utilisation}
\subsection{Cas « S'authentifier »}
L'utilisateur saisit son identifiant et son mot de passe. Le système vérifie ces
informations auprès de la table de sécurité dédiée ; en cas de succès, une
session est ouverte et l'utilisateur est redirigé vers le tableau de bord
correspondant à son rôle ; sinon, un message d'erreur est affiché.

\subsection{Cas « Soumettre / valider une demande »}
Le directeur commercial soumet une demande (création d'un client, d'un employé ou
modification de données), enregistrée avec le statut \texttt{EN\_ATTENTE} et
notifiée à l'administrateur. Ce dernier peut l'approuver ou la refuser en y
associant un commentaire ; le statut évolue alors vers \texttt{EFFECTUE} ou
\texttt{REFUSE}, et l'action est journalisée.

\section*{Conclusion}
Les besoins étant spécifiés et illustrés, le chapitre suivant présente la
conception de l'application et de son volet décisionnel.
